%%%%%%%%%%%%%%%%%%%%%%%%%%%%%%%%%%%%%%%%%%%%%%%%%%%%%%%%%%%%%%%%
\section{文献综述}
%%%%%%%%%%%%%%%%%%%%%%%%%%%%%%%%%%%%%%%%%%%%%%%%%%%%%%%%%%%%%%%%

本研究连接三条文献：LLM风险决策、人类替他人决策，以及金融代理中的模型安全对齐。

\subsection{LLM风险决策与行为偏差}

前景理论指出，人类在收益域通常风险规避，在损失域可能风险寻求，并会对概率进行非线性加权~\cite{Kahneman1979Prospect,Tversky1974Judgment}。
近期研究开始将这一框架用于评估LLM。
Zhang等人发现部分LLM可以拟合累积前景理论中的价值函数和权重函数~\cite{Zhang2025Using}；Payne进一步显示风险框架和语义场景会显著改变模型选择~\cite{Payne2025Analysis}。
但这些规律并不稳定。
Hintze等人发现LLM的概率转换不总是呈现人类典型的逆S形模式~\cite{Hintze2025Nonlinear}；Xiao与Wang也发现GPT系列在不同社会场景中的风险反应可能与前景理论预测相反~\cite{Xiao2025Evaluating}。
因此，LLM风险偏好应被视为可测量但高度依赖模型与提示语境的行为特征。

LLM作为合成被试的可行性已得到初步支持。
Aher等人和Horton等人分别展示了LLM复现实验被试行为和模拟经济代理人的潜力~\cite{Aher2023Using,Horton2024Large}；Brand等人和Bail则讨论了生成式AI在市场研究、调查和社会科学实验中的方法价值与风险~\cite{Brand2023Using,Bail2024Can}。
与此同时，人格提示、社会期望偏差和隐含推理偏差研究表明，角色设定会系统性影响模型输出~\cite{Salecha2024LLMs,Jiang2023Evaluating,Gupta2024Bias,Bini2025Behavioral}。
这为本文通过Prompt操纵“为自己/替他人/是否可观察”提供了方法基础，也提示我们必须谨慎处理提示词本身的诱导效应。
与既有研究相比，本文关注的不是LLM能否在单次回答中给出“像人类”的选择，而是同一模型在受控实验条件下的输出分布是否随角色与问责线索发生系统移动。
这种分布视角更适合评价API模型，因为实际应用中的代理系统往往会在相似任务上被反复调用，单次回答的偶然性并不足以代表其部署风险。

\subsection{人类替他人决策}

人类替他人做风险选择时是否更冒险，一直没有单一答案。
Polman与Wu的元分析纳入71篇研究、128个效应量，发现总体上存在较小的风险寻求偏移，但研究间异质性很高~\cite{Polman2020Decision}。
Fullbrunn与Luhan发现有限责任制度可能降低损失厌恶并提高风险承担，而无奖金或责任更强的条件又可能使代理人更谨慎~\cite{Fullbrunn2017Decision,Fullbrunn2019Responsibility}。
Buckle等人、Pahlke等人进一步说明，责任对象、金额和制度安排会改变代理风险偏好的方向~\cite{Buckle2024Needs,Pahlke2015Responsibility}。

机制层面，后悔、责任感、社会价值判断和认知投入都会影响替他人决策。
Stone等人发现后悔顾虑会增加风险规避，Stone与Allgaier则指出只有在风险被社会认可的领域中，替他人决策才更可能冒险~\cite{Stone2002Risk,Stone2008Social}。
Rigoli等人和Barrafrem与Hausfeld的过程证据显示，为他人决策可能伴随更高随机性或更低信息加工投入~\cite{Rigoli2018Risk,Barrafrem2020Tracing}。
Batteux等人和Michael等人则提示，代理决策可能更接近期望值中性，或主要受社会信息而非结果归属本身影响~\cite{Batteux2017Risk,Batteux2020Decisions,Michael2021Decision}。
这些发现共同说明，代理风险偏移的方向应由情境和问责结构决定，而不是预设为单一效应。
但现有代理决策研究主要基于人类被试，其机制往往依赖真实损益体验、情绪反应、社会评价和制度责任。
LLM没有直接财富感受，也不会以人类方式承担法律或声誉后果；它对“责任”“客户”“损失”等词语的反应更可能来自训练语料、对齐目标和提示语义。
因此，人类代理决策结论不能直接外推到LLM，需要在相同风险量表下重新检验。

\subsection{金融代理与安全对齐}

LLM已被用于金融预测、投资分析和交易代理。
Lopez-Lira与Tang、Fatouros等人和Yu等人分别展示了LLM在股票收益预测、金融推理和交易代理中的应用潜力~\cite{LopezLira2023Can,Fatouros2024Can,Yu2023FinMem}；Saha等人综述了LLM投资代理的基础、基准和挑战~\cite{Saha2025Large}。
但金融场景中的模型偏差同样明显。
Lee等人发现LLM投资分析存在行业、规模和策略偏好，Dimino等人则揭示了金融二分决策中的位置偏差~\cite{Lee2025Your,Dimino2025Tracing}。

安全对齐可能进一步改变代理风险行为。
RLHF和Constitutional AI旨在提升有用性并降低有害输出~\cite{Ouyang2022Training,Bai2022Constitutional}，但后续研究显示，对齐机制会受到微调、语义操纵、人格提示和指令微调动态的影响~\cite{Qi2024Finetuning,McIntosh2024Inadequacy,Liu2024Evaluating,Wu2024Language}。
这意味着金融LLM代理的风险问题不能只从预测准确率或收益率评估。
一个模型即使能够生成合理的市场分析，也可能在“替客户决策”“客户是否能看到过程”“面对收益或损失”这类社会性线索下改变其风险取向。
对于金融应用而言，这类偏移既可能表现为过度保守，导致机会成本上升，也可能表现为低问责场景下的风险放松，放大委托人的潜在损失。
因此，有必要把LLM视为可实验测量的行为主体，对其在代理角色中的风险偏好进行独立审计。
现有研究尚未系统检验：当LLM被设定为金融代理人时，代理角色、问责线索和损益框架是否会共同改变其风险偏好。
本文的实验正是针对这一交叉空白。
